\begin{textsample}{NEG Dim 2 – 2022\_c4 – Score 18.00 – 2022\_c4/enade\_2022\_tec\_log\_15f82e2ab1.txt}  \label{ex:f2_neg_407}
CAPÍTULO 5 PERCEPÇÃO DA PROVA Nas análises feitas neste capítulo, constam as percepções dos concluintes da Área de Tecnologia em Logística sobre a prova aplicada no Enade/2022.
Essas percepções foram mensuradas por meio de nove questões que avaliaram desde o grau de dificuldade da prova até o tempo gasto para concluí -la.
As percepções sobre a prova foram relacionadas ao desempenho dos estudantes, à grande região de funcionamento do curso, à categoria administrativa e ao tipo de organização acadêmica da IES.
O questionário de percepção da prova encontra -se ao final do Anexo VIII, no qual está a reprodução do exame.
Na apresentação dos dados relativos às nove questões sobre as percepções a respeito da prova, o desempenho dos estudantes foi classificado em quatro quartos.
Para tanto, esse desempenho foi ordenado de forma ascendente.
O percentil 25, P25, também conhecido como primeiro quartil, é a nota de desempenho que deixa um quarto ( 25 \% ) dos valores observados abaixo e três quartos acima.
O quarto inferior de desempenho é composto pelas notas abaixo do primeiro quartil.
Já o percentil 75, P75, também conhecido como terceiro quartil, é o valor para o qual há três quartos ( 75 \% ) dos dados abaixo e um quarto acima dele.
O quarto superior de desempenho é composto pelas notas iguais ou acima do terceiro quartil.
O percentil 50, P50, também conhecido como mediana, é o valor que divide as notas em dois conjuntos de igual tamanho.
O segundo quarto inclui os valores entre o primeiro quartil ( P25 ) e a mediana.
O terceiro quarto contém os valores entre a mediana ( P50 ) e o terceiro quartil ( P75 ).
Vale ressaltar que percentis, quartis e medianas são pontos que, não obrigatoriamente, pertencem ao conjunto original de dados, ao passo que os quartos são subconjuntos dos dados originais.
Na Figura 5.1, apresenta -se uma ilustração dos quatro quartos descritos.
A seguir, serão apresentados gráficos com resultados selecionados, relativos às nove questões avaliadas por grupos de estudantes.
Nas barras dos gráficos, apresenta -se o percentual de estudantes que assinalaram uma das opções ou a soma das porcentagens daqueles que assinalaram duas ( ou três ) delas.
Por exemplo, para as questões 1 e 2, nos gráficos, é apresentada a porcentagem total de participantes que assinalaram as opções ( D ) Difícil e ( E ) Muito difícil.
Em cada barra, foram assinalados, também, os extremos do intervalo de confiança de 95 \% como linhas verticais unidas por uma linha horizontal na forma da letra H maiúscula.
O estimador de um parâmetro com um certo nível de confiança ( por exemplo, 95,0 \% ) deve conter o parâmetro no intervalo de confiança em 95 \% das vezes.
Na comparação entre os estimadores dos parâmetros de duas classes de uma dada categoria ( por exemplo : Norte e Nordeste nas grandes regiões, de primeiro e último quarto dentro de desempenho, de IES públicas e privadas, ou de universidades e faculdades ), associados aos seus respectivos intervalos de confiança, diz -se que não há diferença estatisticamente significativa entre os parâmetros das duas categorias quando há interseção entre os intervalos de confiança, e que há diferença se os intervalos de confiança são disjuntos ( para maiores informações, vide Glossário ).
Nas Tabelas do Anexo V, são apresentados os valores absolutos e a distribuição percentual8 das alternativas válidas das nove questões, segundo o mesmo recorte de desempenho dos estudantes e da grande região de funcionamento do curso, de categoria administrativa e do tipo de organização acadêmica da IES. 5.1 GRAU DE DIFICULDADE DA PROVA Nesta seção são mostradas as percepções dos estudantes quanto ao grau de dificuldade da prova Enade/2022, e que foram analisadas considerando sua vinculação ao Componente de Formação Geral e ao Componente Específico. 5.1.1 COMPONENTE DE FORMAÇÃO GERAL O Gráfico 5.1 apresenta as respostas obtidas para a Questão 1 : “ Qual o grau de dificuldade desta prova na parte de Formação Geral? ”.
De acordo com os dados do Gráfico 5.1, a Região Sul apresentou o maior percentual de estudantes que consideraram a prova “ Difícil ” ou “ Muito difícil ”, com 39,4 \%, seguida da região Norte com 37,9 \%.
Na região Nordeste a média foi de 36,2 \%, na Sudeste 33,4 \% e na região Centro-Oeste 35,2 \%.
A região Sudeste apresentou a menor proporção de estudantes que consideraram a prova “ Difícil ” ou “ Muito Difícil ”.
Não há diferenças significativas entre as regiões, exceto entre Sul e Sudeste.
O Gráfico 5.2 traz mais resultados sobre a questão tratada no gráfico anterior, porém, por quartos de desempenho.
Os dados do Gráfico 5.2 mostram que o percentual de estudantes que consideraram o componente de Formação Geral da prova “ Difícil ” ou “ Muito difícil ” foi decrescente em função dos quartos de desempenho : 42,7 \%, no 1º quarto e 28,6 \%, no 4º quarto, o grupo de melhor desempenho na prova.
Nos quartos de desempenho intermediários, a proporção de estudantes que consideraram a prova “ Difícil ” ou “ Muito Difícil ” foi de 38 \%, no 2º quarto, e de 36 \% no 3º quarto.
O Gráfico 5.3 traz outros resultados também em relação à mesma questão retratada no gráfico anterior, porém por organização acadêmica das IES.
Como mostram os dados do Gráfico 5.3, o percentual de estudantes que consideraram o Componente de Formação Geral da prova “ Difícil ” ou “ Muito difícil ” foi maior entre os estudantes de universidades, alcançando 38,3 \%, enquanto entre os estudantes dos CEFET/IF, a organização acadêmica com a menor incidência, o percentual alcançou 26,6 \%.
O Gráfico 5.4 apresenta os dados por categoria administrativa das IES.
Os dados do Gráfico 5.4 mostram que o percentual de estudantes que consideraram o Componente de Formação Geral da prova “ Difícil ” ou “ Muito Difícil ” foi maior entre os estudantes de cursos de IES privadas ( 37,6 \% ), enquanto a proporção registrada em IES públicas foi de 24,9 \%. 5.1.2 COMPONENTE ESPECÍFICO O Gráfico 5.5 apresenta as respostas obtidas para a Questão 2 : “ Qual o grau de dificuldade desta prova na parte de Componente Específico? ”.
De acordo com os dados do Gráfico 5.5, a região Norte apresentou o maior percentual de estudantes que consideraram o Componente Específico da prova “ Difícil ” ou “ Muito Difícil ”, com 45,4 \%, seguida da região Sul com 45,1 \%.
Na região Nordeste a média foi de 42,3 \%, na Centro-Oeste, 41,3 \% e na região Sudeste 39,5 \%, com menor percentual.
Não há diferenças significativas entre as regiões Norte, Nordeste, Centro-Oeste e Sul.
O Gráfico 5.6 traz mais resultados sobre a questão tratada no gráfico anterior, porém, por quartos de desempenho.
Os dados do Gráfico 5.6 mostram que não há grande variação no percentual de estudantes que consideraram o Componente Específico da prova “ Difícil ” ou “ Muito Difícil ” em função dos quartos de desempenho : 41,7 \%, no 1º quarto, e 42,7 \%, no 4º quarto, o grupo de melhor desempenho na prova.
Nos quartos de desempenho intermediários, a proporção de estudantes que consideraram a prova “ Difícil ” ou “ Muito Difícil ” foi de 41,9 \%, no 2º quarto, e de 42,8 \% no 3º quarto.
O Gráfico 5.7 traz outros resultados também em relação à mesma questão retratada no gráfico anterior, porém, por organização acadêmica das IES.
Como mostram os dados do Gráfico 5.7, o percentual de estudantes que consideraram o Componente Específico de Formação Geral da prova “ Difícil ” ou “ Muito Difícil ” foi maior entre os estudantes de centros universitários, alcançando 43,3 \%, enquanto entre os estudantes das faculdades, a organização acadêmica com a menor incidência, o percentual alcançou 37,3 \%.
O Gráfico 5.8 apresenta os dados por categoria administrativa das IES.
Os dados do Gráfico 5.8 mostram que o percentual de estudantes que consideraram o Componente Específico da prova “ Difícil ” ou “ Muito Difícil ” foi maior entre os estudantes de cursos de IES privadas ( 43 \% ), enquanto que a proporção registrada em IES públicas foi de 36,4 \%. 5.2 EXTENSÃO DA PROVA EM RELAÇÃO AO TEMPO TOTAL O Gráfico 5.9 apresenta as respostas obtidas para a Questão 3 : “ Considerando a extensão da prova, em relação ao tempo total, você considera que a prova foi longa ou muito longa? ”.
Em relação aos dados do Gráfico 5.9, entre as grandes regiões, os que avaliaram a prova como longa ou muito longa, em relação ao tempo total destinado à sua resolução, atingiram uma porcentagem que variou de 25 \%, na região Norte, até 30,5 \%, na região Nordeste.
Não há diferenças significativas entre as regiões Nordeste, Centro-Oeste e Sudeste.
O Gráfico 5.10 traz mais resultados sobre a questão tratada no gráfico anterior, porém por quartos de desempenho.
No Gráfico 5.10, observa -se que a proporção de estudantes que consideraram a prova longa ou muito longa foi menor no 4º quarto de desempenho ( 24,3 \% ), e maior, no 1º quarto ( 33,4 \% ), indicando um decréscimo de estudantes que consideraram a prova longa, com o aumento do desempenho.
O Gráfico 5.11 traz outros resultados também em relação à mesma questão retratada no gráfico anterior, porém, por organização acadêmica das IES.
De acordo com os dados do Gráfico 5.11, considerando -se a extensão da prova, o percentual de estudantes que avaliaram a prova como longa ou muito longa, em relação ao tempo total destinado à sua resolução, foi maior entre os estudantes de CEFET/IF, alcançando 30,7 \% e com variação maior que os demais, enquanto entre os estudantes de universidades, a organização acadêmica com a menor incidência, esse índice atingiu 28 \%.
Entre os participantes de centros universitários, a proporção atingiu 29,7 \%, e entre os de faculdades, 28,4 \%.
O Gráfico 5.12 apresenta os dados por categoria administrativa das IES.
Os dados do Gráfico 5.12 mostram que o percentual de estudantes que consideraram a prova longa ou muito longa, em relação ao tempo total destinado à sua resolução, foi ligeiramente maior entre os estudantes de cursos de IES privadas ( 28,8 \% ) do que entre os estudantes de IES públicas ( 28,3 \% ).
Não há diferenças significativas entre os dados das duas categorias. 5.3 COMPREENSÃO DOS ENUNCIADOS DAS QUESTÕES Também organizados em relação ao Componente de Formação Geral e ao Componente Específico, os dados apresentados nesta seção dizem respeito à forma como os estudantes compreenderam os enunciados de questões da prova, considerando as formulações expressas nesses mesmos enunciados. 5.3.1 COMPONENTE DE FORMAÇÃO GERAL O Gráfico 5.13 apresenta as respostas obtidas para a Questão 4 : “ Os enunciados das questões da prova na parte de Formação Geral estavam claros e objetivos? ”.
Os resultados apresentados no Gráfico 5.13 mostram que, na análise regional, a porcentagem de estudantes que avaliaram que todos ou a maioria dos enunciados das questões do Componente de Formação Geral estavam “ claros e objetivos ” variou de 73,8 \%, na região Norte até 77,8 \%, na região Sudeste.
Não há diferenças significativas entre as regiões.
O Gráfico 5.14 traz mais resultados sobre a questão tratada no gráfico anterior, porém por quartos de desempenho.
De acordo com os dados do Gráfico 5.14, observa -se que a proporção dos que concordaram com a afirmativa cresce, conforme o desempenho aumenta.
No quarto superior, a clareza e a objetividade de todos ou da maioria dos enunciados das questões foram percebidas por 84,6 \% dos estudantes.
O Gráfico 5.15 traz outros resultados também em relação à mesma questão retratada no gráfico anterior, porém por organização acadêmica das IES.
Os dados do Gráfico 5.15 mostram que a proporção dos estudantes que consideraram que todos ou a maioria dos enunciados das questões do Componente de Formação Geral estavam “ claros e objetivos ” foi maior entre aqueles que estudam em CEFET/IF, chegando a 79,2 \%, enquanto entre os que estudam em universidades, a organização acadêmica com a menor incidência, esse índice atingiu 74,7 \%.
Entre os estudantes de centros universitários, o índice atingiu 76,9 \%, e entre os de faculdades, 79 \%.
Não há diferenças significativas entre centros universitários, faculdades e CEFET/IF.
O Gráfico 5.16 apresenta os dados por categoria administrativa das IES.
Os dados do Gráfico 5.16 mostram que o percentual de estudantes que consideraram que todos ou a maioria dos enunciados das questões do Componente de Formação Geral estavam “ claros e objetivos ” foi maior entre os estudantes de cursos de IES públicas ( 81,2 \% ) quando comparado aos estudantes de IES privadas ( 75,5 \% ). 5.3.2 COMPONENTE ESPECÍFICO O Gráfico 5.17 apresenta as respostas obtidas para a Questão 5 : “ Os enunciados das questões da prova na parte de Componente Específico estavam claros e objetivos? ”.
A maioria dos estudantes de todas as grandes regiões brasileiras considerou “ claros e objetivos ” todos ou a maioria dos enunciados das questões do Componente Específico da prova, percentual sempre maior ou igual a 72,2 \%, a menor proporção.
A maior proporção foi de 75,7 \%, na região Sudeste.
Não há diferenças significativas entre as regiões.
O Gráfico 5.18 traz mais resultados sobre a questão tratada no gráfico anterior, porém, por quartos de desempenho.
Os dados do Gráfico 5.18 mostram que a proporção de estudantes que consideraram os enunciados das questões “ claros e objetivos ” apresenta uma tendência crescente em relação ao aumento de desempenho : mais elevada no quarto superior ( 82,8 \% ), se comparada ao quarto inferior de desempenho ( 64,1 \% ).
O Gráfico 5.19 traz outros resultados, também, em relação à mesma questão retratada no gráfico anterior, porém, por organização acadêmica das IES.
Os dados do Gráfico 5.19 mostram que a proporção dos presentes que consideraram que todos ou a maioria dos enunciados das questões do Componente Específico estavam “ claros e objetivos ” foi maior entre os estudantes dos CEFET/IF e das faculdades, atingindo 79,3 \%, cada.
Entre os estudantes de universidades e de centros universitários, ambas organizações acadêmicas com as menores incidências, essa proporção atingiu 73,3 \% e 73,5 \%, respectivamente.
Não há diferenças significativas entre CEFET/IF e faculdades e entre universidades e centros universitários ou entre centros universitários e CEFET/IF.
O Gráfico 5.20 apresenta os dados por categoria administrativa das IES.
Os dados do Gráfico 5.20 mostram que o percentual de estudantes que consideraram que todos ou a maioria dos enunciados das questões do Componente Específico estavam “ claros e objetivos ” foi maior entre os estudantes de cursos de IES públicas ( 80,9 \% ) que entre os de IES privadas ( 73,5 \% ). 5.4 SUFICIÊNCIA DAS INFORMAÇÕES/INSTRUÇÕES FORNECIDAS O Gráfico 5.21 apresenta as respostas obtidas para a Questão 6 : “ As informações/instruções fornecidas para a resolução das questões foram suficientes para resolvê -las? ”.
De acordo com os dados do Gráfico 5.21 e quanto à distribuição de respondentes pelas grandes regiões, observa -se que a proporção de estudantes que consideraram as informações/instruções fornecidas até excessivas ou suficientes em todas ou na maioria das questões foi sempre superior ou igual a 79,1 \% ( proporção mais baixa na região Sul ), chegando a 82,8 \%, na região Sudeste, maior índice.
Há diferença significativa entre as regiões Sul e Sudeste.
O Gráfico 5.22 traz resultados sobre a questão tratada no gráfico anterior, porém, por quartos de desempenho.
De acordo com os dados do Gráfico 5.22, levando -se em conta o desempenho dos participantes, observa -se uma tendência crescente das proporções de participantes que avaliaram as informações/instruções como até excessivas ou suficientes em todas ou na maioria das questões.
O percentual foi mais elevado no quarto superior ( 87,3 \% ).
Já no quarto inferior, a suficiência das informações/instruções declarada como até excessiva, em todas ou na maioria das questões, foi percebida por 73,7 \% dos respondentes.
O Gráfico 5.23 traz outros resultados também em relação à mesma questão retratada no gráfico anterior, porém, por organização acadêmica das IES.
De acordo com os dados do Gráfico 5.23, considerando -se as informações/instruções fornecidas para a resolução das questões, a maior proporção dos que consideraram que essas eram até excessivas ou suficientes em todas ou na maioria das questões ocorreu entre os estudantes de faculdades, totalizando 85,3 \%.
As universidades foram a organização acadêmica em que houve a menor proporção : 79,6 \%.
Não há diferenças significativas entre universidades e centros universitários e entre faculdades e CEFET/IF e há diferenças significativas entre universidades e faculdades.
O Gráfico 5.24 apresenta os dados por categoria administrativa das IES.
De acordo com os dados do Gráfico 5.24, o percentual de estudantes que consideraram que as informações/instruções fornecidas para a resolução das questões eram até excessivas ou suficientes em todas ou na maioria das questões foi maior entre os estudantes de cursos de IES públicas ( 87,2 \% ), enquanto a proporção registrada entre os de IES privadas foi de 80 \%.
Há diferença significativa entre as duas categorias. 5.5 DIFICULDADE ENCONTRADA AO RESPONDER À PROVA O Gráfico 5.25 apresenta as respostas obtidas para a Questão 7 : “ Você se deparou com alguma dificuldade ao responder à prova?
Qual? ”.
Nos Gráficos 5.25 a 5.28, apresentados de forma sequenciada mais adiante, são expostos os percentuais de estudantes que apontaram o “ Desconhecimento do conteúdo ” como “ dificuldade ao responder à prova ”.
O Gráfico 5.25 apresenta as respostas obtidas sobre o “ Desconhecimento do conteúdo ” como a principal “ dificuldade ao responder à prova ”.
De acordo com os dados do Gráfico 5.25, o percentual de estudantes que apontaram o desconhecimento do conteúdo como a principal “ dificuldade ao responder à prova ” não superou 21,3 \%, maior índice na região Nordeste.
O menor percentual, 17,5 \%, deu -se na região Centro-Oeste.
Há diferença significativa entre as regiões Sudeste e Nordeste.
O Gráfico 5.26 traz resultados sobre o ponto tratado no gráfico anterior, porém, por quartos de desempenho.
De acordo com os dados do Gráfico 5.26, em relação aos quartos de desempenho, o desconhecimento do conteúdo foi a opção escolhida por 18,1 \% dos estudantes do quarto inferior e por 19,4 \% do quarto superior.
O segundo quarto apresentou uma proporção de 19,3 \%, maior que a do terceiro quarto, que foi de 18,2 \%.
O Gráfico 5.27 traz outros resultados também em relação à mesma questão retratada no gráfico anterior, porém, por organização acadêmica das IES.
De acordo com os dados do Gráfico 5.27, na análise por organização acadêmica, o percentual de inscritos e presentes que apontaram o Desconhecimento do conteúdo como “ dificuldade ao responder à prova ” não superou 16,8 \%.
Os percentuais variaram de 16,8 \%, nas faculdades, até 23,1 \%, nos CEFET/IF.
Não há diferenças significativas entre as organizações acadêmicas.
O Gráfico 5.28 apresenta os dados por categoria administrativa das IES.
De acordo com os dados do Gráfico 5.28, em relação às categorias administrativas, o desconhecimento do conteúdo foi a opção escolhida por 19,1 \% dos estudantes nas IES privadas e 16,2 \%, nas IES públicas.
Não há diferença significativa entre IES públicas e privadas, considerando a intersecção dos intervalos de confiança no valor 18,3 \% 5.6 CONTEÚDOS DAS QUESTÕES OBJETIVAS DA PROVA O Gráfico 5.29 apresenta as respostas obtidas para a Questão 8 : “ Considerando apenas as questões objetivas da prova, você percebeu que não estudou ainda a maioria desses conteúdos? ”.
Na análise por grande região, a proporção de respondentes que escolheram a opção não estudou ainda a maioria desses conteúdos foi pequena.
Observa -se que, nas regiões Norte ( 7,8 \% ) e Nordeste ( 10,3 \% ), as proporções foram ligeiramente maiores.
Na região Sudeste, a situação foi inversa, com proporções mais baixas, correspondendo a 6,4 \%.
Há diferenças significativas das regiões Sul e Sudeste com a região Nordeste.
Não há diferenças significativas entre as regiões Norte, Nordeste e Centro-Oeste.
O Gráfico 5.30 traz resultados sobre a questão tratada no gráfico anterior, porém, por quartos de desempenho.
De acordo com os dados do Gráfico 5.30, considerando -se separadamente as opiniões dos estudantes dos quatro quartos de desempenho, observa -se que, no quarto inferior, 13,2 \% ofereceram como resposta que não estudaram, ainda, a maioria desses conteúdos.
A menor proporção foi dos estudantes do quarto superior, correspondendo a 2,8 \%.
Levando -se em conta o desempenho dos participantes, observa -se uma tendência decrescente das proporções de participantes que informaram não ter estudado ainda a maioria dos conteúdos, do primeiro ao último quarto.
O Gráfico 5.31 traz outros resultados também em relação à mesma questão retratada no gráfico anterior, porém, por organização acadêmica das IES.
De acordo com os dados do Gráfico 5.31, na análise por organização acadêmica, a proporção de respondentes que escolheram a opção não estudou ainda a maioria desses conteúdos foi pequena.
Observa -se que a proporção nas faculdades ( 4,9 \% ) foi menor.
Nos centros universitários, a situação foi inversa : a proporção foi mais alta, 8,8 \%.
As IES CEFET/IF e universidades apresentaram valores intermediários, sendo 6,9 \% e 7,5 \% respectivamente.
Há diferenças significativas das universidades e centros universitários em relação às faculdades.
O Gráfico 5.32 apresenta os dados por categoria administrativa das IES.
De acordo com os dados do Gráfico 5.32, em relação às categorias administrativas, observa -se que, nas IES públicas, 3,2 \% ofereceram como resposta que não estudou, ainda, a maioria desses conteúdos, sendo 8,1 \% nas IES privadas com a mesma resposta.
Há diferença significativa entre as categorias administrativas. 5.7 TEMPO GASTO PARA CONCLUIR A PROVA O Gráfico 5.33 apresenta as respostas obtidas para a Questão 9 : “ Qual foi o tempo gasto por você para concluir a prova? ”.
Considerando -se as cinco grandes regiões brasileiras, os percentuais dos que utilizaram entre duas e quatro horas para finalizar a prova foram, nas regiões Norte e Nordeste, 73,6 \% e 73,1 \%, respectivamente, sendo percentuais muito próximos entre essas regiões.
Nas regiões Sul e Centro-Oeste, os percentuais de estudantes que dispensaram entre duas e quatro horas para concluir a prova foram 67,8 \% e 70 \%, respectivamente.
A região Sudeste apresentou proporção de 69,8 \%.
Há diferenças significativas entre a região Sul e as regiões Norte e Nordeste.
O Gráfico 5.34 traz resultados sobre a questão tratada no gráfico anterior, porém, por quartos de desempenho.
De acordo com os dados do Gráfico 5.34, uma vez analisadas as alternativas escolhidas pelos estudantes que se situam nos diferentes quartos de desempenho, observa -se uma tendência crescente : uma proporção maior de participantes no quarto superior ( 79,6 \% ) declarou ter gasto entre duas e quatro horas para concluir a prova quando comparada com a proporção de participantes nos quartos inferiores.
A proporção no quarto inferior foi de 58,1 \%.
O Gráfico 5.35 traz outros resultados também em relação à mesma questão retratada no gráfico anterior, porém por organização acadêmica das IES.
De acordo com os dados do Gráfico 5.35, na análise por organização acadêmica, os percentuais dos que utilizaram de duas a quatro horas para finalizar a prova foram de 71,5 \% nos centros universitários ; 72,2 \% nas faculdades e 78,5 \% nos CEFET/IF.
O menor percentual foi entre os estudantes de universidades, com 68,4 \% dos alunos necessitando de duas a quatro horas para concluir a prova.
Há diferenças significativas entre as universidades e as demais organizações acadêmicas.
O Gráfico 5.36 apresenta os dados por categoria administrativa das IES.
De acordo com os dados do Gráfico 5.36, em relação às categorias administrativas, observa -se que nas IES públicas, 73 \% declararam ter gasto entre duas e quatro horas para concluir a prova.
Registra -se 69,9 \% nas IES privadas com a mesma resposta.
Há diferença significativa entre IES públicas e privadas.
Os resultados do Enade 2022 apresentados neste capítulo contribuem para o entendimento de que a consideração da percepção dos estudantes concluintes sobre a prova do Enade é importante tanto para a compreensão do seu desempenho como para a reflexão sistemática e constante a respeito de determinadas realidades relacionadas a esse mesmo desempenho no país.
Por outro lado, como todos os elementos do trabalho pedagógico institucional mostram -se diretamente relacionados com o desempenho dos estudantes, também ganha relevância conhecer e analisar os resultados da percepção dos coordenadores de curso sobre a mesma prova.
As informações expressas por estudantes e coordenadores sobre questões diversas associadas aos dois componentes da prova – Formação Geral e Conhecimento Específico – podem agregar valor a ações institucionais voltadas para o aperfeiçoamento do instrumento avaliativo, considerando os vários elementos envolvidos em sua elaboração.
Neste contexto, é relevante que os resultados das percepções de estudantes e coordenadores sobre a prova do Enade ganhem peso como um elemento a mais nas análises da oferta dos cursos nas modalidades presencial e a distância, em nível nacional.
Espera -se, assim, que, em uma visão formativa de avaliação, o conhecimento e a análise desses resultados possam vir a contribuir para a construção de alternativas institucionais que venham a induzir a qualidade da educação superior.

% matched lemmas: 
\end{textsample}
